\section{Бетатронное движение в магнитном кольце}
\label{sec:ch2}

В главе~\ref{sec:magnetic} период задан интегралом~\eqref{eq:mag-T}. Скорость вынесена из него, и для одной дисперсионной орбиты ряд доведён до $\eta_1$. Здесь в тот же элемент длины возвращаются бетатронное отклонение и углы. Модуль скорости между зазорами по-прежнему зависит только от полного импульса: это свойство статического магнитного поля. Коэффициент, с которым бетатронное удлинение входит в период, получается из произведения $(v_s/v)(C/C_s)$.

\subsection{Элемент пути}
\label{sec:ch2-dl}

Элемент пути в криволинейной координате $s$ есть
\begin{equation}
\label{eq:beta-dl}
\mathrm{d}l
=\sqrt{\Bigl(1+\frac{x}{\rho}\Bigr)^2+x'^2+y'^2}\;\mathrm{d}s.
\end{equation}
Здесь $\rho$ --- локальный радиус кривизны опорной орбиты; вне диполей $1/\rho=0$. При $x/\rho>-1$ множитель $1+x/\rho$ выносится точно,
\begin{equation}
\label{eq:beta-dl-factor}
\mathrm{d}l
=\Bigl(1+\frac{x}{\rho}\Bigr)
\sqrt{1+\frac{x'^2+y'^2}{(1+x/\rho)^2}}\;\mathrm{d}s.
\end{equation}
Квадрат $(x/\rho)^2$ в отдельное слагаемое не выделяется: он уже сидит внутри $1+x/\rho$. Разложение корня по квадрату угла до второго порядка по поперечным амплитудам даёт
\begin{equation}
\label{eq:beta-dl-2}
\mathrm{d}l
=\Biggl(1+\frac{x}{\rho}+\frac{x'^2+y'^2}{2}\Biggr)\mathrm{d}s.
\end{equation}
Произведение $(x/\rho)(x'^2+y'^2)$ имеет третий порядок и отбрасывается.

Отклонение складывается из дисперсионной орбиты, бетатронного колебания и сдвига, квадратичного по бетатронной амплитуде,
\begin{equation}
\label{eq:beta-x}
x=D\delta+x_\beta+x_{\mathrm{nl}},
\qquad
x'=D'\delta+x'_\beta+\cdots,
\qquad
y'=y'_\beta+\cdots.
\end{equation}
Слагаемые, составленные только из $D$, $D'$ и $\delta$, включая $(D'\delta)^2/2$, уже собраны в $\alpha_0$ и $\alpha_1$ формулы~\eqref{eq:mag-C}. Подстановка~\eqref{eq:beta-x} в~\eqref{eq:beta-dl-2} оставляет сверх той длины четыре бетатронных вклада: линейный $x_\beta/\rho$, квадрат угла $(x_\beta'^2+y_\beta'^2)/2$, перекрёстное произведение $D'\delta\,x'_\beta$ и сдвиг $x_{\mathrm{nl}}/\rho$.

\subsection{Среднее по бетатронной фазе}
\label{sec:ch2-avg}

Бетатронный набег фазы за оборот имеет порядок единицы, синхротронный тюн --- порядок $10^{-3}$. За время, на котором синхротронная фаза ещё почти постоянна, бетатронная фаза успевает много раз обойти окружность. Действие $J$ на этом медленном времени есть адиабатический инвариант. Высокочастотное поле видит опоздание одного оборота, а уравнение фазы усредняет его ещё и по бетатронной фазе. У Шоджи~\cite{Shoji2005} тому же пределу отвечает $n\to\infty$ в формуле (5): среднее удлинение одного оборота извлекается из разности пути за большое число оборотов.

Линейный интеграл $\oint x_\beta/\rho\,\mathrm{d}s$ после такого усреднения равен нулю. Сумма по оборотам остаётся ограниченной. В разделе II.B работы~\cite{Shoji2005} амплитуда продольного смещения есть $\sqrt{2J_x\mathcal{H}}$, где $\mathcal{H}=\gamma_x D^2+2\alpha_x DD'+\beta_x D'^2$, формулы (10)--(12). Постоянной части, растущей вместе с числом оборотов, у этого члена нет. Перекрёстное произведение $D'\delta\,x'_\beta$ линейно по бетатронной амплитуде, и его среднее по фазе тоже равно нулю. Квадрат $(x/\rho)^2$ геометрией~\eqref{eq:beta-dl-factor} не порождается, поэтому произведение $x_\beta D\delta$ в длину на этом порядке не входит.

Квадрат угла остаётся. Решение линейного уравнения с параметрами Твисса имеет вид $x_\beta=\sqrt{2J_x\beta_x}\cos\psi$, $\mathrm{d}\psi/\mathrm{d}s=1/\beta_x$ и $\beta_x'=-2\alpha_x$. Дифференцирование по $s$ даёт
\begin{equation}
\label{eq:beta-xp}
x'_\beta
=-\sqrt{\frac{2J_x}{\beta_x}}\,(\alpha_x\cos\psi+\sin\psi).
\end{equation}
На фиксированном азимуте фаза за много оборотов распределена равномерно, если тюн не целый. Средние $\langle\cos^2\psi\rangle=\langle\sin^2\psi\rangle=1/2$ и $\langle\sin\psi\cos\psi\rangle=0$ превращают квадрат~\eqref{eq:beta-xp} в
\begin{equation}
\label{eq:beta-xp2}
\langle x_\beta'^2\rangle
=\frac{2J_x}{\beta_x}\Biggl(\frac{\alpha_x^2}{2}+\frac{1}{2}\Biggr)
=J_x\gamma_x,
\end{equation}
где $\gamma_x=(1+\alpha_x^2)/\beta_x$. Локальное $\alpha_x$ в тождестве сохраняется. То же по вертикали, $\langle y_\beta'^2\rangle=J_y\gamma_y$. Среднее углового интеграла за оборот есть
\begin{equation}
\label{eq:beta-angle}
\Biggl\langle\oint\frac{x_\beta'^2+y_\beta'^2}{2}\,\mathrm{d}s\Biggr\rangle
=\frac{J_x}{2}\int\gamma_x\,\mathrm{d}s
+\frac{J_y}{2}\int\gamma_y\,\mathrm{d}s.
\end{equation}
Функция $\alpha=-\beta'/2$ периодична, так что $\int\alpha\,\mathrm{d}s=0$, но в $\langle x'^2\rangle$ стоит $\alpha^2$, и интеграл $\int(\alpha^2/\beta)\,\mathrm{d}s$ остаётся. Замена $\gamma$ на $1/\beta$ этот интеграл выбрасывает.

На кольце из $16$ ячеек FODO с тонкими линзами локальный $\alpha$ отличен от нуля. Дрейф между линзами равен $0{,}8$~м, сила фокусирующей линзы равна $0{,}4$~м$^{-1}$, дефокусирующая имеет обратный знак. На этом кольце $\int\gamma_x\,\mathrm{d}s=10{,}374$ и $2\pi\nu_x=5{,}142$; вертикальная плоскость совпадает с горизонтальной по симметрии ячейки. Сумма $x'^2\,\mathrm{d}s/2+y'^2\,\mathrm{d}s/2$ по $8000$ оборотам при $J_x=10^{-6}$~м и $J_y=0{,}7\times 10^{-6}$~м относится к правой части~\eqref{eq:beta-angle} как $1{,}00000035$. Та же сумма больше величины, в которой $\gamma$ заменено на $1/\beta$, в $2{,}017$ раза.

Сдвиг $x_{\mathrm{nl}}$ квадратичен по бетатронной амплитуде. Его вклад $\oint x_{\mathrm{nl}}/\rho\,\mathrm{d}s$ имеет тот же порядок, что и~\eqref{eq:beta-angle}, и в среднее входит. Для короткого секступоля он сводится к толчку, усреднённому по бетатронной фазе; кривизна и краевое поле дают свои квадратичные члены. Их перечень --- в приложении~\ref{app:ch2path}. Оба вклада, угловой и нелинейный, не зависят от $\delta$ в удерживаемом порядке: хроматическая поправка $\gamma(\delta)=\gamma+\delta\,\partial\gamma/\partial\delta+\cdots$ умножает действие на $\delta$ и отбрасывается вместе с произведениями $\delta\Lambda$.

Одно относительное удлинение собирает всё, что осталось после усреднения,
\begin{equation}
\label{eq:beta-Lambda}
\Lambda
\equiv
\frac{1}{C_s}
\Biggl(
\frac{J_x}{2}\int\gamma_x\,\mathrm{d}s
+\frac{J_y}{2}\int\gamma_y\,\mathrm{d}s
+\Biggl\langle\oint\frac{x_{\mathrm{nl}}}{\rho}\,\mathrm{d}s\Biggr\rangle
\Biggr).
\end{equation}
Первая группа слагаемых есть прямое содержание квадрата угла. Вторая есть средний сдвиг орбиты. Дальше в период входит именно это $\Lambda$, и второй буквы для того же вклада нет.

\subsection{Множитель скорости}
\label{sec:ch2-speed}

Длина оборота теперь есть сумма дисперсионной длины и бетатронного удлинения,
\begin{equation}
\label{eq:beta-C}
\frac{C}{C_s}
=1+\alpha_0\delta+\alpha_1\delta^2+\Lambda.
\end{equation}
Слагаемые $\alpha_0\delta$ и $\alpha_1\delta^2$ содержат только $D$ и $D'$. Величина $\Lambda$ прибавлена интегралом~\eqref{eq:beta-dl-2} внутри $C$, до умножения на $v_s/v$. Поскольку $|v|$ между зазорами зависит только от $\delta$, формула~\eqref{eq:mag-Tratio} остаётся точной:
\begin{equation}
\label{eq:beta-Tratio}
\frac{T-T_s}{T_s}
=\frac{v_s}{v}\,\frac{C}{C_s}-1.
\end{equation}
Подстановка ряда~\eqref{eq:mag-vsv} и суммы~\eqref{eq:beta-C} повторяет произведение~\eqref{eq:mag-product} и добавляет к нему $(v_s/v)\Lambda$. При
\begin{equation}
\label{eq:beta-vsv-L}
\frac{v_s}{v}\,\Lambda
=\Lambda-\frac{\Lambda}{\gamma_s^2}\,\delta+O(\delta^2\Lambda)
\end{equation}
удержанный член есть само $\Lambda$. Следующий, $-(\Lambda/\gamma_s^2)\delta$, вместе с прочими произведениями $\delta\Lambda$ отбрасывается. Поэтому
\begin{equation}
\label{eq:beta-T}
\frac{T-T_s}{T_s}
=\eta_0\delta+\eta_1\delta^2+\Lambda,
\end{equation}
с теми же $\eta_0$ и $\eta_1$, что в~\eqref{eq:mag-eta0} и~\eqref{eq:mag-eta1}.

Коэффициент перед $\Lambda$ равен единице, потому что $v_s/v$ при $\delta=0$ равен единице, а $\Lambda$ уже второго порядка малости. Коэффициент $\alpha_0$ стоит при $\delta$ внутри дисперсионной части $C/C_s$ и на $\Lambda$ не умножается: вклады в интеграл~\eqref{eq:beta-dl-2} складываются. Запись удлинения множителем $(1+\lambda)$ при дисперсионной длине, $\lambda=\Delta C_\beta/C_{\mathrm{disp}}$, даёт в произведении $(1+\eta_0\delta)\lambda$ слагаемое $\eta_0\delta\lambda$. Пересчёт $\lambda=\Lambda\,C_s/C_{\mathrm{disp}}=\Lambda(1-\alpha_0\delta+\cdots)$ снимает $\alpha_0\delta\Lambda$, потому что $\eta_0-\alpha_0=-1/\gamma_s^2$, и первая поправка снова равна $-(\delta/\gamma_s^2)\Lambda$. В удержанном порядке перед $\Lambda$ по-прежнему стоит единица. Величина $\eta_0=\alpha_0-1/\gamma_s^2$ есть коэффициент при $\delta$ в~\eqref{eq:beta-T} и знаменатель сдвига импульса; в множитель при $\Lambda$ она на этом порядке не входит.

Порядок отброшенного произведения виден на числах магнитного кольца. Амплитуда синхротронных колебаний $\delta_m$ лежит между $10^{-4}$ и $10^{-3}$. Действие при эмиттансе $\varepsilon=2J=10^{-6}$~м равно $J=5\times 10^{-7}$~м. Ниже, в записи через хроматичность, этому действию в обеих плоскостях отвечает $\Lambda=9{,}875\times 10^{-8}$. Произведение $\delta\Lambda$ тогда лежит между $10^{-11}$ и $10^{-10}$: при $\delta_m=10^{-4}$ член $|\Lambda|\delta_m/\gamma_s^2$ равен $7{,}5\times 10^{-12}$. Рядом с ним удержанные $\Lambda$ и $\eta_1\delta_m^2/2=5{,}1\times 10^{-9}$ больше на три--четыре порядка. Поправка, которую $-(\Lambda/\gamma_s^2)\delta$ вносит в $\eta_0$, имеет величину $\Lambda/\gamma_s^2\sim 8\times 10^{-8}$ и в $\delta_{\mathrm{eq}}$ не входит.

Тот же коэффициент единица стоит в формуле (6) работы~\cite{Senichev2013}, куда $(\Delta L/L)_\beta$ вписано после того, как продольное уравнение уже сокращено с линейной скоростью. Совпадение этого коэффициента с~\eqref{eq:beta-T} следует из интеграла~\eqref{eq:beta-Tratio}. Вклад входит как кусок $T=\oint\mathrm{d}l/v$ после усреднения по быстрому бетатрону, и перед геометрическим удлинением в этом куске стоит единица.

\subsection{Равновесный импульс}
\label{sec:ch2-eq}

Уравнение фазы есть~\eqref{eq:mag-phase-slip} с рядом~\eqref{eq:beta-T} вместо~\eqref{eq:mag-eta-series}. Уравнение энергии на зазоре остаётся~\eqref{eq:mag-energy}. При $\cos\varphi=1$ и $\omega_s^2=-K\omega_{\mathrm{rf}}\eta_0$ из~\eqref{eq:mag-ws2} получается
\begin{equation}
\label{eq:beta-osc}
\ddot\delta+\omega_s^2\delta
=-\frac{\omega_s^2}{\eta_0}(\eta_1\delta^2+\Lambda).
\end{equation}
На синхротронном времени $\Lambda$ постоянно. У периодического решения среднее $\ddot\delta$ равно нулю, а $\langle\delta^2\rangle=\delta_m^2/2$ с точностью до членов четвёртого порядка: постоянный сдвиг сам имеет порядок $\delta_m^2$ и $\Lambda$. Поэтому
\begin{equation}
\label{eq:beta-avg}
\eta_0\langle\delta\rangle+\eta_1\frac{\delta_m^2}{2}+\Lambda=0,
\end{equation}
\begin{equation}
\label{eq:beta-deq}
\delta_{\mathrm{eq}}
=\langle\delta\rangle
=-\frac{1}{\eta_0}
\Biggl(\frac{\eta_1\delta_m^2}{2}+\Lambda\Biggr).
\end{equation}
Это магнитный аналог формулы (10) в~\cite{Senichev2013} с полным $\eta_1$ и с тем знаком, который уже стоит в среднем уравнения (8) той работы: раздел~\ref{sec:sign} показал, что печатный предымножитель $\gamma_s^2/(\gamma_s^2\alpha_0-1)=1/\eta_0$ этот минус теряет. При $\Lambda=0$ формула~\eqref{eq:beta-deq} возвращается к~\eqref{eq:mag-deq}.

Решение до второй гармоники даёт приложение~\ref{app:ch1bm}. Для~\eqref{eq:beta-osc} постоянная $c_0=-(\omega_s^2/\eta_0)\Lambda$ и $c_2=-(\omega_s^2/\eta_0)\eta_1$, так что
\begin{equation}
\label{eq:beta-delta-sol}
\delta(t)
=\delta_m\cos\psi
-\frac{1}{\eta_0}
\Biggl(\frac{\eta_1\delta_m^2}{2}+\Lambda\Biggr)
+\frac{\eta_1\delta_m^2}{6\eta_0}\cos 2\psi.
\end{equation}
Постоянная часть совпадает с~\eqref{eq:beta-deq}. Частота на том же порядке есть формула~\eqref{eq:mag-bm-B2-result},
\begin{equation}
\label{eq:beta-freq}
\frac{\omega}{\omega_s}
=1-\frac{5}{12}\frac{\eta_1^2\delta_m^2}{\eta_0^2}
-\frac{\eta_1\Lambda}{\eta_0^2}.
\end{equation}
Подстановка~\eqref{eq:beta-deq} в~\eqref{eq:beta-T} обращает среднее опоздание в нуль,
\begin{equation}
\label{eq:beta-sync}
\Biggl\langle\frac{T-T_s}{T_s}\Biggr\rangle
=\eta_0\delta_{\mathrm{eq}}+\eta_1\frac{\delta_m^2}{2}+\Lambda
=0.
\end{equation}
Сдвиг $\delta_{\mathrm{eq}}$ исчезает при любой амплитуде $\delta_m$ и при любых действиях, когда оба коэффициента обращаются в нуль,
\begin{equation}
\label{eq:beta-comp}
\eta_1=0,
\qquad
\Lambda=0.
\end{equation}
В записи раздела~\ref{sec:ch2-xi} второе равенство для коротких секступолей есть $\xi_x=\xi_y=0$. Три коэффициента $\eta_1$, $\xi_x$ и $\xi_y$ независимы, и сил секступолей, которые их гасят, нужно три семейства. Этот счёт прочитан из~\eqref{eq:beta-sync}. Интеграл по кольцу собирает только суммы вида $\int\beta DK_2\,\mathrm{d}s$ и не знает порядка расстановки элементов; однооборотное отображение при той же оптике может дать другой сдвиг.

При $\eta_0<0$ и $\Lambda>0$ формула~\eqref{eq:beta-deq} даёт $\delta_{\mathrm{eq}}>0$. Положительный сдвиг импульса укорачивает период и компенсирует бетатронное удлинение так же, как он компенсирует положительное $\eta_1$ в конце раздела~\ref{sec:equilibrium}.

\subsection{Гладкое кольцо и формула (11)}
\label{sec:ch2-smooth}

В гладком кольце $\alpha=0$, поэтому $\gamma=1/\beta$, и $\int\mathrm{d}s/\beta=2\pi\nu$, поскольку $\mathrm{d}\psi=\mathrm{d}s/\beta$. Угловая часть~\eqref{eq:beta-Lambda} без $x_{\mathrm{nl}}$ принимает вид
\begin{equation}
\label{eq:beta-smooth}
\Lambda
=\frac{J_x\nu_x+J_y\nu_y}{2R}
=\frac{\pi(J_x\nu_x+J_y\nu_y)}{C_s},
\end{equation}
где $C_s=2\pi R$. Эмиттанс формулы (11) в~\cite{Senichev2013} определён так, что фазовое среднее $x_\beta'^2=\varepsilon_x/(2\beta_x)$. Вместе с $\alpha=0$ это есть $\varepsilon=2J$: из~\eqref{eq:beta-xp2} при $\alpha=0$ следует $\langle x'^2\rangle=J/\beta=\varepsilon/(2\beta)$. Подстановка $J=\varepsilon/2$ в~\eqref{eq:beta-smooth} даёт
\begin{equation}
\label{eq:beta-wepea11}
\Lambda
=\frac{\pi}{2C_s}(\varepsilon_x\nu_x+\varepsilon_y\nu_y).
\end{equation}
Печатная формула (11) той работы есть $(\Delta L/L)_\beta=\pi(\varepsilon_x\nu_x+\varepsilon_y\nu_y)/(2L)$ при $L=C_s$. Промежуточная запись там же есть $(1/L)\int\bigl(x_\beta/\rho+(x_\beta'^2+y_\beta'^2)/2\bigr)\mathrm{d}s$ при $x_\beta'^2=\varepsilon/(2\beta)$ и $1/\beta=\nu/R$. Произведение двух множителей $1/2$, геометрического и того, что сидит в $\varepsilon/(2\beta)$, даёт
\begin{equation}
\label{eq:beta-wepea-alg}
\frac{1}{L}\int\frac{1}{2}\cdot\frac{\varepsilon}{2\beta}\,\mathrm{d}s
=\frac{\varepsilon}{4L}\int\frac{\mathrm{d}s}{\beta}
=\frac{\varepsilon}{4L}\cdot 2\pi\nu
=\frac{\pi\varepsilon\nu}{2L}.
\end{equation}
Печатный множитель $\pi/(2L)$ с этим промежуточным счётом совпадает. Расхождения в двойке между формулой (11) и строкой $x'^2=\varepsilon/(2\beta)$ нет. Двойка появилась бы, если выбросить геометрическую половину в~\eqref{eq:beta-dl-2} или подставить в~\eqref{eq:beta-wepea11} эмиттанс $\varepsilon_{\mathrm{rms}}=\langle J\rangle$ вместо $\varepsilon=2J$. В формуле (20) той же работы этот $\varepsilon$ назван среднеквадратичным; при определении $\langle J\rangle=\varepsilon_{\mathrm{rms}}$ такая подстановка занижает $\Lambda$ вдвое. Для пучка среднее~\eqref{eq:beta-Lambda} линейно по $J$, и в него входит $\langle J\rangle=\varepsilon_{\mathrm{rms}}$.

В неопубликованной заметке~\cite{SenichevNote} шаг перед формулой (9) читается так: среднее $\beta'$ по кольцу равно нулю, следовательно $\alpha=0$, и тогда $x'_\beta=-\sqrt{2J/\beta}\,\sin\psi$ совпадает с $-\sqrt{2J\gamma}\,\sin\psi$. Среднее $\int\beta'\,\mathrm{d}s=0$ верно, потому что $\beta$ периодична. Из него не следует $\alpha(s)=0$ на каждом азимуте. Равенство $\sqrt{1/\beta}=\sqrt{\gamma}$ при $\alpha\neq 0$ ложно. Если в интеграл подставить амплитуду $\sqrt{2J/\beta}$ и отбросить $\alpha\cos\psi$, получится $(J/2)\int\mathrm{d}s/\beta$ и пропадёт $(J/2)\int(\alpha^2/\beta)\,\mathrm{d}s$. Формула (11) заметки записана уже через $\gamma$ и равна $(J/2)\int\gamma\,\mathrm{d}s$ на каждую плоскость. Этот интеграл совпадает с~\eqref{eq:beta-angle}: его спасает тождество $\langle x'^2\rangle=J\gamma$, доказанное в~\eqref{eq:beta-xp2} без предположения $\alpha=0$. Ложный шаг --- отождествление двух амплитуд в формуле (9). Конечная формула (11) заметки с~\eqref{eq:beta-angle} согласована.

\subsection{Хроматичность}
\label{sec:ch2-xi}

Вторая запись того же $\Lambda$ выражает нелинейный сдвиг через хроматичность. В модели второго порядка работы~\cite{Shoji2005}, с секступолем, кривизной и краевым полем, формулы (26) и (30) после интегрирования по частям сокращаются до формулы (31),
\begin{equation}
\label{eq:beta-xi}
\Lambda
=-\frac{2\pi}{C_s}(J_x\xi_x+J_y\xi_y).
\end{equation}
Здесь $\xi=\mathrm{d}\nu/\mathrm{d}\delta$. Абсолютное удлинение Шоджи есть $\Delta C=-2\pi(J_x\xi_x+J_y\xi_y)$, и~\eqref{eq:beta-xi} есть оно же, делённое на $C_s$. Эта запись равна~\eqref{eq:beta-Lambda}, потому что хроматичность собрана из тех же сил, что и путь: квадрат угла даёт интеграл $\gamma$, а секступоль, кривизна и краевое поле дают $\langle\oint x_{\mathrm{nl}}/\rho\,\mathrm{d}s\rangle$.

Интеграл одной $\gamma$ совпадает с~\eqref{eq:beta-xi}, когда среднее $\oint x_{\mathrm{nl}}/\rho\,\mathrm{d}s$ равно нулю. Для этого должны отсутствовать секступольный вклад, вклад кривизны сверх того, что уже сидит в функции $\gamma(s)$ линейной оптики, и вклад краевого поля. Тогда $\xi_x=-(1/(4\pi))\int\gamma_x\,\mathrm{d}s$ и так же по вертикали, с фокусирующим коэффициентом уравнения $z''+Kz=0$, и обе записи сводятся к~\eqref{eq:beta-angle}, делённому на $C_s$. Пока эти вклады есть, основная запись полного $\Lambda$ --- формула~\eqref{eq:beta-xi}. Запись через $\gamma$ остаётся прямой оценкой квадрата угла и становится полным $\Lambda$, когда $x_{\mathrm{nl}}$ в среднем не удлиняет орбиту.

Короткий секступоль --- условие, при котором заметка~\cite{SenichevNote} заменяет $x_{\mathrm{nl}}$ средним толчком и приходит к своей формуле (24), бетатронная часть которой совпадает с~\eqref{eq:beta-xi}. Кривизна и краевое поле в этот вывод не входят. У Шоджи они удерживаются и в пути, и в хроматичности, и равенство (31) при них сохраняется. Обнуление $\Lambda$ вместе с $\xi_x$ и $\xi_y$ в приближении коротких секступолей есть частный случай того же равенства. В модели Шоджи $\Lambda$ обнуляется вместе с полной хроматичностью, уже включающей кривизну и край.

Тождество, которое переводит квадруполь в $\gamma$, относится к коэффициенту $K$ в уравнении $z''+Kz=0$:
\begin{equation}
\label{eq:beta-twiss}
\int K\beta\,\mathrm{d}s=\int\gamma\,\mathrm{d}s,
\end{equation}
поскольку $\alpha'=K\beta-\gamma$ и $\int\alpha'\,\mathrm{d}s=0$. В горизонтальной плоскости $K_x=K_1+1/\rho^2$, в вертикальной $K_y=-K_1$, если $K_1$ --- градиент, входящий в $x''+K_1 x$ со знаком плюс. Формула (22) заметки пишет $\int K_1\beta_{x,y}\,\mathrm{d}s=\int\gamma_{x,y}\,\mathrm{d}s$ для обеих плоскостей. По горизонтали в ней нет $1/\rho^2$. По вертикали знак $K_1$ противоположен $K_y$.

Формула (21) заметки имеет вид
\begin{equation}
\label{eq:beta-note21}
\xi_x=-\frac{1}{4\pi}\int(K_1\beta_x+K_2\beta_x D)\,\mathrm{d}s,
\qquad
\xi_y=-\frac{1}{4\pi}\int(K_1\beta_y-K_2\beta_y D)\,\mathrm{d}s.
\end{equation}
Горизонтальная строка согласована с уравнением $x''+K_1 x=(K_2/2)(x^2-y^2)+\cdots$ и с толчком, который даёт путь; в ней нет дипольного $1/\rho^2$. Вертикальная строка ставит $K_1$ с тем же знаком, что и горизонтальная. Уравнение вертикали есть $y''-K_1 y=-K_2 xy$, поэтому фокусирующий коэффициент есть $-K_1$. В~\eqref{eq:beta-note21} перед $K_1\beta_y$ стоит тот же минус, что и в горизонтальной строке. Знак $K_2$ в вертикальной строке с тем же толчком согласован: обе плоскости дают в длине $(J_x/2)\int K_2\beta_x D\,\mathrm{d}s-(J_y/2)\int K_2\beta_y D\,\mathrm{d}s$.

Подстановка ложного тождества $\int K_1\beta_y=\int\gamma_y$ в ложный знак $K_1$ возвращает $\xi_y=-(1/(4\pi))\int(\gamma_y-K_2\beta_y D)\,\mathrm{d}s$. Это то выражение, которое переводит угловой и секступольный интегралы формулы (20) заметки в $-2\pi J_y\xi_y$. Две ошибки в квадрупольном куске компенсируют друг друга, и бетатронная часть формулы (24) с путём (20) сходится. Дипольный член $1/\rho^2$ отсутствует и в (21), и в (22), поэтому его нет и в (24). Естественная горизонтальная хроматичность этот член содержит.

Формула (32) в~\cite{Shoji2005} закрывает удлинение сдвигом импульса
\begin{equation}
\label{eq:beta-shoji32-main}
\delta=-\frac{1}{\alpha}\frac{\Delta C}{C_0}.
\end{equation} Делитель $\alpha$ там есть коэффициент уплотнения. Для электронов на $1$~ГэВ величина $1/\gamma^2$ мала рядом с обычным $\alpha$; в квазиизохронной строке таблицы~I той работы $\alpha=2\times 10^{-6}$ уже сравним с $1/\gamma^2$. Для дейтронов магнитного кольца $\alpha_0=0{,}230$ и $\eta_0=-0{,}534211$ различаются и знаком. Подстановка $\alpha_0$ вместо $\eta_0$ в~\eqref{eq:beta-deq} меняет знак $\delta_{\mathrm{eq}}$. Равновесие по-прежнему читается из~\eqref{eq:beta-sync}, и в знаменателе стоит $\eta_0$.

Формула равновесного сдвига в статье об оптимизации спиновой когерентности методом отображения записана как
\begin{equation}
\label{eq:beta-map}
\Delta\delta_{\mathrm{eq}}
=\frac{\gamma^2}{\gamma^2\alpha_0-1}
\Biggl[
\frac{1}{2}\,\delta_m^2\,\eta_1
-\frac{\pi}{L}(\varepsilon_x\xi_x+\varepsilon_y\xi_y)
\Biggr].
\end{equation}
Предымножитель равен $1/\eta_0$, как в~\eqref{eq:mag-prefactor}. Коэффициент $\eta_1$ в той формуле совпадает с~\eqref{eq:mag-eta1}, включая $3\beta_s^2/(2\gamma_s^2)$. При $\varepsilon=2J$ и $L=C_s$ бетатронное слагаемое скобки есть
\begin{equation}
\label{eq:beta-map-L}
-\frac{\pi}{L}(\varepsilon_x\xi_x+\varepsilon_y\xi_y)
=-\frac{2\pi}{C_s}(J_x\xi_x+J_y\xi_y)
=\Lambda.
\end{equation}
Скобка~\eqref{eq:beta-map} совпадает со скобкой~\eqref{eq:beta-deq}. Предымножитель $1/\eta_0$ имеет знак, противоположный~\eqref{eq:beta-deq}: это та же потеря минуса, что в печатных (10) и (16) работы~\cite{Senichev2013}. Относительный знак члена $\eta_1\delta_m^2/2$ и бетатронного члена при этом согласован: оба описывают удлинение пути. Если в~\eqref{eq:beta-map} подставить $\varepsilon_{\mathrm{rms}}=\langle J\rangle$, множитель $\pi/L$ даёт половину $\Lambda$. Среднее по пучку в~\eqref{eq:beta-xi} есть замена $J$ на $\langle J\rangle=\varepsilon_{\mathrm{rms}}$.

\subsection{Числа магнитного кольца}
\label{sec:ch2-num}

Длина $C_s$ есть сумма элементов решётки \texttt{magnetic\_2.fox}: шестнадцать диполей по $5{,}3294148$~м, четыре прямых промежутка той же длины, квадруполи, секступоли и дрейфы. Сумма равна $117{,}388$~м. Таблица OptiM заканчивается на $11738{,}83$~см. В паспорте кольца $\xi_x=-1{,}84$, $\xi_y=-1{,}85$, $\eta_0=-0{,}534211$, $\eta_1=1{,}02$. В списке OptiM силы секступолей равны нулю, так что это естественная хроматичность, и~\eqref{eq:beta-xi} содержит кривизну и краевое поле наравне с квадратом угла.

Для частицы с $\varepsilon_x=\varepsilon_y=10^{-6}$~м в смысле $\varepsilon=2J$, то есть $J_x=J_y=5\times 10^{-7}$~м,
\begin{equation}
\label{eq:beta-num-L}
\Lambda
=-\frac{2\pi}{C_s}\bigl(J_x\xi_x+J_y\xi_y\bigr)
=9{,}875\times 10^{-8}.
\end{equation}
При $\delta_m=10^{-4}$ квадратичный член импульса равен $\eta_1\delta_m^2/2=5{,}10\times 10^{-9}$. Бетатронное слагаемое больше этого члена в $19{,}4$ раза. Сдвиг~\eqref{eq:beta-deq} равен $1{,}944\times 10^{-7}$: из него $1{,}849\times 10^{-7}$ даёт $\Lambda$ и $9{,}547\times 10^{-9}$ даёт $\eta_1$. При $\delta_m=10^{-3}$ член $\eta_1\delta_m^2/2=5{,}10\times 10^{-7}$ больше $\Lambda$ в $5{,}16$ раза, и $\delta_{\mathrm{eq}}=1{,}140\times 10^{-6}$.

Если эмиттанс пучка задан как $\varepsilon_{\mathrm{rms}}=\langle J\rangle=10^{-6}$~м в каждой плоскости, действие в~\eqref{eq:beta-num-L} удваивается и $\Lambda=1{,}975\times 10^{-7}$. Гладкая формула~\eqref{eq:beta-smooth} на тюнах таблицы OptiM той же решётки, $\nu_x=2{,}462$ и $\nu_y=2{,}461$, даёт при $J=5\times 10^{-7}$~м значение $6{,}59\times 10^{-8}$. Она заменяет $\gamma$ на $\nu/R$ и не содержит $x_{\mathrm{nl}}$; с~\eqref{eq:beta-num-L} она совпадать не обязана.
